\documentclass{beamer}

\input{MathMacros}

\usepackage{beamerthemesplit}
\usepackage{enumerate}
\usepackage{amssymb}
\usepackage[all,cmtip]{xy}
\usepackage[mathscr]{eucal}
\usepackage[natbib=true,style=authoryear,backend=bibtex,useprefix=true]{biblatex}
\addbibresource{reading.bib}
\usepackage{graphicx,color}

\title{Writing and Thinking with LLMs: Discipline Over Convenience}
\author{Reuben Brasher}
\date{\today}

\begin{document}

\frame{\titlepage}

\section[Outline]{}
\frame{\tableofcontents}

\section{Lecture 2}

\frame
{
   \frametitle{Writing Is a Commodity, Thinking Is Not}

    Writing tasks are now easy to automate, but thinking remains scarce.

    \begin{itemize}
        \item Large language models can generate essays, summaries, and reports on demand.
        \item The barrier to producing polished text has effectively disappeared.
        \item Superficial fluency is no longer a meaningful indicator of skill.
        \item Audiences are increasingly familiar with AI-generated patterns.
        \item Low-quality outputs are recognizable despite surface polish.
        \item Differentiation has shifted from producing text to producing insight.
    \end{itemize}
}

\frame
{
   \frametitle{The Collapse of Naive Writing Tasks}

    Traditional writing tasks no longer reliably measure competence.

    \begin{itemize}
        \item One-page essays, memos, and summaries can be generated instantly.
        \item Standard corporate formats are easily replicated by LLMs.
        \item These outputs have become commoditized across organizations.
        \item Readers can detect generic or formulaic AI-generated content.
        \item The presence of text no longer implies understanding or effort.
        \item Evaluation shifts from whether it was written to whether it is meaningful.
    \end{itemize}
}

\frame
{
   \frametitle{The Corporate Reality of Communication}

    Effective writing in professional contexts depends on clarity and credibility.

    \begin{itemize}
        \item Most stakeholders do not read long documents in detail.
        \item Careful readers evaluate quality immediately.
        \item Generic or verbose content undermines credibility.
        \item AI-generated patterns are increasingly recognizable in business settings.
        \item Influence depends on clarity of thought rather than quantity of text.
        \item Writing must support real-time communication and decision-making.
    \end{itemize}
}

\frame
{
   \frametitle{The New Role of the Human}

    The human role shifts from generating text to directing and evaluating it.

    \begin{itemize}
        \item Humans must define the problem and interpret results.
        \item Real-time thinking remains essential in meetings and presentations.
        \item Arguments must be formed independently of the tool.
        \item LLMs cannot replace dynamic, context-sensitive responses.
        \item The user is responsible for coherence and correctness.
        \item The tool amplifies capability but does not replace judgment.
    \end{itemize}
}

\frame
{
   \frametitle{Analogy: Thinking Without Assistance}

    Human cognition must remain functional without constant tool support.

    \begin{itemize}
        \item Tools assist but are not always available in real-time situations.
        \item Presentations and discussions require immediate reasoning.
        \item Internal competence enables effective use of external tools.
        \item Reliance without understanding leads to failure under pressure.
        \item Mastery requires both tool use and independent capability.
        \item The tool is an extension rather than a substitute for cognition.
    \end{itemize}
}

\frame
{
   \frametitle{Grounded Writing: Start with Sources}

    Effective writing begins with grounding in real information.

    \begin{itemize}
        \item Conduct a literature or document search relevant to the topic.
        \item Review abstracts and summaries to assess relevance.
        \item Select high-quality academic, industry, or news sources.
        \item Avoid beginning with open-ended prompts to the model.
        \item Use external information to constrain the model’s output.
        \item Establish factual grounding before generating text.
    \end{itemize}
}

\frame
{
   \frametitle{Grounded Writing: Input Real Content}

    The model should operate on real documents rather than free-form prompts.

    \begin{itemize}
        \item Provide PDFs, articles, or curated text as input.
        \item Extract key elements such as titles, authors, and sections.
        \item Request structured summaries of each source.
        \item Ensure the model operates within provided material.
        \item Reduce reliance on the model’s internal distribution.
        \item Treat the model as a processor of information, not a source.
    \end{itemize}
}

\frame
{
   \frametitle{Grounded Writing: Build Structure}

    Structured summaries create a stable foundation for writing.

    \begin{itemize}
        \item Compile an annotated bibliography from selected sources.
        \item Generate concise summaries for each reference.
        \item Maintain consistent citation formats.
        \item Use structured outputs to organize information.
        \item Create reusable reference sets for future work.
        \item Build a foundation before drafting content.
    \end{itemize}
}

\frame
{
   \frametitle{Controlling the Model: Explicit Instruction}

    Clear instruction is required to constrain model behavior.

    \begin{itemize}
        \item Instruct the model to avoid meta text and unnecessary commentary.
        \item Specify stylistic constraints such as tone and voice.
        \item Prohibit unwanted formatting unless explicitly needed.
        \item Define structure and output expectations clearly.
        \item Reinforce constraints repeatedly when necessary.
        \item Treat instruction as an ongoing process.
    \end{itemize}
}

\frame
{
   \frametitle{Controlling the Model: Active Correction}

    Users must intervene actively to maintain correctness and relevance.

    \begin{itemize}
        \item Identify and call out hallucinations immediately.
        \item Reject outputs that deviate from the task.
        \item Reinforce boundaries when the model drifts.
        \item Do not accept partially correct responses.
        \item Guide the model back to the intended objective.
        \item Maintain control over the conversation direction.
    \end{itemize}
}

\frame
{
   \frametitle{Controlling the Model: Multiple Conversations}

    Separating tasks into distinct contexts improves reliability.

    \begin{itemize}
        \item Use one conversation for content generation.
        \item Use another for editing and critique.
        \item Assign distinct roles to each interaction.
        \item Prevent context contamination between tasks.
        \item Maintain focused and manageable interactions.
        \item Treat conversations as components of a workflow.
    \end{itemize}
}

\frame
{
   \frametitle{Question-Driven Writing Method}

    Writing improves when used to structure thinking rather than replace it.

    \begin{itemize}
        \item Define the topic and provide relevant sources.
        \item Ask the model to generate structured questions.
        \item Use these questions to guide your thinking.
        \item Answer each question independently.
        \item Build content from your own responses.
        \item Ensure ideas originate from the human.
    \end{itemize}
}

\frame
{
   \frametitle{Incremental Drafting}

    Documents should be built in controlled sections rather than generated at once.

    \begin{itemize}
        \item Draft one section at a time.
        \item Provide context and constraints for each section.
        \item Avoid generating full documents in a single prompt.
        \item Maintain control over structure and progression.
        \item Reduce drift and inconsistency.
        \item Assemble sections into a coherent whole.
    \end{itemize}
}

\frame
{
   \frametitle{Iterative Refinement: Improving Quality}

    High-quality writing emerges through repeated evaluation and revision.

    \begin{itemize}
        \item Remove redundant content across sections.
        \item Ensure consistent tone and voice.
        \item Eliminate generic or formulaic phrasing.
        \item Improve clarity and precision of language.
        \item Review transitions between sections.
        \item Perform multiple refinement passes.
    \end{itemize}
}

\frame
{
   \frametitle{Iterative Refinement: Multi-Model Review}

    Using multiple models can improve detection of weaknesses.

    \begin{itemize}
        \item Use a second model for critique and analysis.
        \item Identify patterns typical of AI-generated text.
        \item Focus on sentence-level improvements.
        \item Detect verbosity and repetition.
        \item Apply targeted refinements.
        \item Combine strengths of different systems.
    \end{itemize}
}

\frame
{
   \frametitle{Ethical Boundary: Ownership and Attribution}

    LLM use introduces questions about authorship and responsibility.

    \begin{itemize}
        \item Fully generated content presented as original is plagiarism.
        \item Editing self-authored content is acceptable.
        \item Intermediate cases require careful judgment.
        \item Ownership depends on origin of ideas.
        \item Proper citation may be necessary.
        \item Ethical use requires transparency.
    \end{itemize}
}

\frame
{
   \frametitle{Discussion}

    Writing with LLMs requires discipline and structured control.

    \begin{itemize}
        \item LLMs are powerful but unreliable without constraints.
        \item Quality depends on grounding in real information.
        \item Structure and workflow determine reliability.
        \item Human oversight is required at every stage.
        \item Effective use depends on discipline, not convenience.
        \item Advantage comes from directing and constraining the tool.
    \end{itemize}
}

\begin{frame}[t,allowframebreaks]
\frametitle{References}
\printbibliography
\end{frame}

\end{document}